\section*{第 \ref{ch:b2:basins-compared} 章——沉积盆地比较}

\begin{solution}{exo:b2:basins-compared:1}
(a) \omterm{def:b2:basins-compared:types}{克拉通内盆地}：沉降缓慢，呈脉动式进行，历时数亿年，无明显
断层作用；其成因尚有争议。(b) 裂谷：拉伸期间地壳经历力学减薄，
随后岩石圈冷却，发生热收缩。(c) 被动大陆边缘：热沉降
继承自裂谷阶段，由厚层沉积楔的负荷放大。(d) 与 (e)
为\omterm{def:b2:alpine-orogen:foreland}{前陆盆地}：岩石圈受阿尔卑斯山与喜马拉雅山逆冲带的负荷而向下挠曲。
\end{solution}

\begin{solution}{exo:b2:basins-compared:2}
(a) $Y = 3000\times(3330 - 2300)/(3330 - 1030) = 3000\times1030/2300 = \qty{1343}{m}$.
(b) $1343 + 200 = \qty{1543}{m}$。(c) 若无此负荷，基底将位于 \qty{1343}{m}：
$1 - 1343/3000 = \qty{55}{\percent}$ 的厚度由负荷腾出空间。
\end{solution}

\begin{solution}{exo:b2:basins-compared:3}
(a) $\beta = 35/25 = 1.4$. (b) $125/1.4 = \qty{89.3}{km}$. (c) $S_i = 4334\times(1 - 1/1.4) =
\qty{1238}{m}$.
\end{solution}

\begin{solution}{exo:b2:basins-compared:4}
(a) $\tau = (1.25\times10^{5})^2/(\pi^2\times8\times10^{-7}) = \qty{1.979e15}{s}$，即
$1.979\times10^{15}/3.156\times10^{13} = \qty{62.7}{Ma}$. (b) $62.7\times0.8 = \qty{50.2}{Ma}$.
(c) $62.7\times(100/125)^2 = \qty{40.1}{Ma}$. (d) $1 - e^{-1} = \qty{63}{\percent}$; $1 -
e^{-3} = \qty{95}{\percent}$.
\end{solution}

\begin{solution}{exo:b2:basins-compared:5}
(a) $0.63\,e^{-1000/1960} = 0.378$; $0.63\,e^{-2000/1960} = 0.227$. (b) $T_s = 1000 -
0.63\times1960\times(0.6004 - 0.3605) = \qty{704}{m}$. (c) $m = 2700\times704 + 1030\times296 =
\qty{2.206e6}{kg.m^{-2}}$; $\bar\rho = \qty{2206}{kg.m^{-3}}$.
\end{solution}

\begin{solution}{exo:b2:basins-compared:6}
(a) $z_2' \leftarrow 704 + 1234.8\,(1 - e^{-z_2'/1960})$: 1077, 1226, 1278, 1296, 1301, 1303,
1304；该层厚 \qty{1304}{m}。(b) $(2700\times704 + 1030\times600)/1304 =
\qty{1932}{kg.m^{-3}}$。(c) 压实挤出的水
量相当于 \qty{304}{m}，即 \qty{30}{\percent}。
\end{solution}

\begin{solution}{exo:b2:basins-compared:7}
(a) $1 - e^{-100/62.7} = 0.7971$。当 $\beta = 1.5$ 时：$3207\times0.4135\times0.7971 =
\qty{1057}{m}$；当 $\beta = 2$ 时：$3207\times0.6366\times0.7971 = \qty{1627}{m}$。(b) 是的，1100 至
\qty{1400}{m} 位于两者之间。(c) $\frac{\beta}{\pi}\sin\frac{\pi}{\beta} = 1250/(3207\times0.7971)
= 0.4890$，故 $\beta = 1.6 + 0.1\times(0.4890 - 0.4705)/(0.5204 - 0.4705) = 1.64$，位于该盆地推断拉伸量的
50 至 \qty{100}{\percent} 范围之内。
\end{solution}

\begin{solution}{exo:b2:basins-compared:8}
(a) 充填物的密度取代分母中的 $\rho_w$。(b) 分子为
$125\,000\times(1/3)\times74.69 = \num{3.112e6}$; $S_i = 3.112\times10^{6}/(3184.4 - 2300) =
\qty{3519}{m}$。(c) 若位于水下，$3.112\times10^{6}/2154.4 = \qty{1445}{m}$：充满沉积物的
盆地深度是其 2.44 倍。
\end{solution}

\begin{solution}{exo:b2:basins-compared:9}
(a) $2171 - 100\times3330/2300 = 2171 - 145 = \qty{2026}{m}$。(b) $Y$ 随
$W_d$ 等量变化：由 1971 增至 \qty{2371}{m}。(c) 水深 $\pm\qty{200}{m}$，
与海平面的 \qty{145}{m} 相比。
\end{solution}

\begin{solution}{exo:b2:basins-compared:10}
(a) $S_i$ 的符号与其分子中括号的符号相同，因为分母 $3184.4 - 1030$ 为
正；括号为负时则为负，即隆升。(b) $(\rho_m - \rho_c)t_c/a =
\alpha T_1\rho_m/2$，得 $t_c = a\,\alpha T_1\rho_m/(2(\rho_m - \rho_c)) = 125\times0.04372\times
3330/1060 = \qty{17.17}{km}$。(c) $t_c = 17.17/(1 - 0.04372\times17.17/250) = \qty{17.22}{km}$。
(d) 拉伸前地壳薄于约 \qty{17}{km} 的地区：如大洋岩石圈，
或已被早期裂谷作用减薄的大陆地壳中。
\end{solution}

\begin{solution}{exo:b2:basins-compared:11}
(a) 拉伸前，该柱体中的岩石圈质量为 $M = \rho_c t_c(1 - \alpha T_1
t_c/(2a)) + \rho_m(a - t_c)(1 - \alpha T_1(a + t_c)/(2a))$；拉伸后，$\rho_w S_i + M/\beta +
\rho_m(1 - \alpha T_1)(a - S_i - a/\beta)$。令两者相等，$S_i[\rho_m(1 - \alpha T_1) - \rho_w] =
(1 - 1/\beta)[\rho_m(1 - \alpha T_1)a - M]$。展开 $M$ 后，$\rho_m(1 - \alpha T_1)a - M =
(\rho_m - \rho_c)t_c(1 - \alpha T_1 t_c/(2a)) - \rho_m\alpha T_1 a/2$：即所求公式。(b) 末态剖面为 $T_1 z/a$；初态剖面比其
多出 $T_1 z(\beta - 1)/a$（$a/\beta$ 以上），以下为 $T_1(1 - z/a)$。积分分别为 $T_1
a(\beta - 1)/(2\beta^2)$ 与 $T_1 a(\beta - 1)^2/(2\beta^2)$，二者之和为 $T_1 a(1 - 1/\beta)/2$。
收缩量即此值乘以 $\alpha$，经水负载后得该表达式。(c) $3.28\times10^{-5}
\times3330\times1333\times125\,000\times0.5/(2\times2300) = \qty{1978}{m}$，对比 $3207\times
2/\pi = \qty{2042}{m}$：第一项高估了末态沉降 \qty{64}{m}，
相对误差为 \qty{3}{\percent}。
\end{solution}

\begin{solution}{exo:b2:basins-compared:12}
(a) $3207\times0.6366\times(1 - e^{-150/62.7}) = 2042\times0.9086 = \qty{1855}{m}$. (b) $1855\times
2300/(3330 - 2400) = \qty{4588}{m}$。(c) 6.9 至 \qty{13.8}{km}。(d) 该模型忽略了若干
可能加深盆地的因素：地幔岩石圈的拉伸程度超过地壳，区域
（挠曲）对沉积负荷的响应，裂谷期间边缘的瞬变隆升
随后下沉，或沉积物供给过多，使盆地过度充填，
并在陆架上堆积至海平面以上。
\end{solution}

\begin{solution}{pb:b2:basins-compared:1}
\textbf{1.} 顶部 $0.49\,e^{-4420/3704} = 0.149$；底部 $0.49\,e^{-6030/3704} = 0.096$。
\textbf{2.} $T_s = 1610 - 0.49\times3704\times(0.3032 - 0.1963) = \qty{1416}{m}$.
\textbf{3.} 页岩：$1110 - 0.63\times1960\times(e^{-3310/1960} - e^{-4420/1960}) = 1110 - 1234.8
\times(0.1847 - 0.1049) = \qty{1011}{m}$.
\textbf{4.} 泥质砂岩：$2680\times1004 + 1030\times696 = \qty{3.408e6}{kg.m^{-2}}$；白垩：
$2710\times1409 + 1030\times201 = \qty{4.025e6}{kg.m^{-2}}$；页岩：$2700\times1011 + 1030\times99 =
\qty{2.832e6}{kg.m^{-2}}$；沙：$2650\times1416 + 1030\times194 = \qty{3.952e6}{kg.m^{-2}}$。合计
\qty{1.4217e7}{kg.m^{-2}}，$\bar\rho_s = 14\,216\,910/6030 = \qty{2358}{kg.m^{-3}}$。
\textbf{5.} $z_2' \leftarrow 1416 + 1815\,(1 - e^{-z_2'/3704})$: 1993, 2171, 2221, 2235, 2238, 2239,
2239: \qty{2239}{m}.
\textbf{6.} $(2650\times1416 + 1030\times823)/2239 = \qty{2055}{kg.m^{-3}}$.
\textbf{7.} $Y = 2239\times(3330 - 2055)/2300 + 200 = 1241 + 200 = \qty{1441}{m}$.
\textbf{8.} $Y = 6030\times(3330 - 2358)/2300 + 100 = 2548 + 100 = \qty{2648}{m}$。各岩层
按现今深度计量：它们的压实程度已与今天相同。
\textbf{9.} $(2171 - 1441)/50 = \qty{14.6}{m}$ 每百万年，自 150 至 \qty{100}{Ma}；
$(2453 - 2171)/40 = \qty{7.05}{m}$，自 100 至 \qty{60}{Ma}；$(2648 - 2453)/60 = \qty{3.25}{m}$
此后沉降减慢。
\textbf{10.} $6030 + 100 = \qty{6130}{m}$，为构造沉降的 2.31 倍。
\textbf{11.} $(6130 - 2648)/6130 = \qty{57}{\percent}$.
\textbf{12.} $2171 - 100\times3330/2300 = \qty{2026}{m}$.
\textbf{13.} 裂谷结束时的点，即 \qty{150}{Ma}：$S_i = \qty{1441}{m}$，前提是
裂谷的 \qty{15}{Ma} 期间未发生热沉降。
\textbf{14.} $\beta = 1/(1 - 1441/4334) = 1/0.6675 = 1.50$.
\textbf{15.} 地壳 $35/1.50 = \qty{23.3}{km}$，岩石圈 $125/1.50 = \qty{83.3}{km}$。
\textbf{16.} $(1.50/\pi)\sin(\pi/1.50) = 0.4775\times0.8660 = 0.4135$; $3207\times0.4135 =
\qty{1326}{m}$.
\textbf{17.} $S_i = 4334/3 = \qty{1445}{m}$。在 \qty{100}{Ma}，$t = \qty{50}{Ma}$：$1326\times(1 -
e^{-50/62.7}) = \qty{729}{m}$，$Y = \qty{2173}{m}$；在 \qty{60}{Ma}：$1326\times0.7620 =
\qty{1010}{m}$，$Y = \qty{2455}{m}$；今天：$1326\times0.9086 = \qty{1205}{m}$，$Y = \qty{2650}{m}$。
\textbf{18.} 模型在各点都深 \qty{2}{m}：2173 对 2171，2455 对 2453，
2650 对 2648。
\textbf{19.} 古水深的估计误差仅能达到百米或更大；过去的海平面
与今天不同；裂谷作用持续 \qty{15}{Ma}，期间已有部分热量
已经散失；采用挠曲而非局地地壳均衡；地幔的拉伸程度多于或少于地壳；
裂谷前地壳厚度并非 \qty{35}{km}。
\textbf{20.} $1326\times(e^{-150/62.7} - e^{-200/62.7}) = \qty{67}{m}$.
\textbf{21.} $66.6\times2300/(3330 - 2400) = \qty{165}{m}$.
\textbf{22.} $1 - e^{-150/62.7} = \qty{91}{\percent}$.
\textbf{23.} 裂谷结束后经过 $\tau\ln2 = \qty{43.5}{Ma}$，即 \qty{106.5}{Ma}。
\textbf{24.} $\beta = \textbf{1.50}$：地壳被拉伸至原宽度的 1.5 倍。
\end{solution}
